\documentclass[11pt]{article}
\usepackage{amsmath,amssymb,amsthm}
\usepackage[margin=1in]{geometry}
\newtheorem{lemma}{Lemma}
\newtheorem{theorem}{Theorem}
\newtheorem{corollary}{Corollary}
\usepackage[hidelinks]{hyperref}
\newcommand{\Q}{\mathbb Q}
\begin{document}
\setlength{\emergencystretch}{1em}
\title{A scalar cyclic minor for ordinary quintic binomial moments}
\author{}\date{}\maketitle
We prove consecutive coprimality at every ordinary order for the
quintic support $\{2,5\}$. A mixed-adjoint identity confines every
common root to one three-adic scale. A scalar minor then becomes
injective on the entire integral disk. Conjugacy and real positivity
exclude a common factor. The existing support reduction consequently
gives Strong Factorial for all ordinary univariate binomials of degree
at most five. Internal replay is distinct from independent replay,
formal verification, external review and priority.
For $n\ge1$ put
\[
 P_n(t)=\sum_{k=0}^n\binom nk(2n+1)_{3k}t^k,\qquad
 D=5n+4,\qquad r=v_3(D).
\]
Every valuation is normalized by $v_3(3)=1$.
We work in a complete algebraically closed extension $K$ of $\Q_3$,
with valuation ring $\mathcal O$, maximal ideal $\mathfrak m$, and residue
field $k$. An overbar always means reduction modulo $\mathfrak m$;
it does not assert congruence modulo $3\mathcal O$.

\section{Two elementary lemmas}
\begin{lemma}[An affine reduction is injective on the integral disk]
\label{lem:affine-injectivity}
Let $F(X)\in\mathbb Z_3[X]$ satisfy $\bar F=c(1+X)$ for some
$c\in\mathbb F_3^\times$. For distinct $x,y\in\mathcal O$,
\[
 v_3(F(x)-F(y))=v_3(x-y).
\]
In particular $F$ has at most one zero in $\mathcal O$, and that
zero is simple.
\end{lemma}
\begin{proof}
Write $F=\sum a_jX^j$. Its linear coefficient is a unit and
$a_j\in3\mathbb Z_3$ for $j\ge2$. The divided difference is
\[
 \frac{F(x)-F(y)}{x-y}
 =a_1+\sum_{j\ge2}a_j\sum_{\ell=0}^{j-1}
                          x^{j-1-\ell}y^\ell.
\]
This is a unit: all terms after the first lie in
$3\mathcal O\subset\mathfrak m$.
The same observation applies to $F'(x)$.
All sums are finite.
\end{proof}

\begin{lemma}[Real-field descent from a unique local candidate]
\label{lem:real-field-descent}
Let $E\subset\mathbb R$ be a field of characteristic zero with an
embedding $\iota:E\longrightarrow K$.
Suppose $P,Q\in E[t]$ have no common real root and their images
under $\iota$ have at most one distinct common root in $K$.
Then $P,Q$ are coprime over $E$ and over $\mathbb C$.
\end{lemma}
\begin{proof}
A nonconstant gcd has an irreducible factor $h$ over $E$.
Characteristic zero makes $h$ separable, so all $\deg h$ distinct
roots of $\iota(h)$ are common roots in $K$.
Thus $\deg h\le1$, giving a root in $E\subset\mathbb R$,
a contradiction. A Bezout identity over $E$ remains one over
$\mathbb C$.
This uses all conjugates of a common factor. It does not claim
that an arbitrary unique integral root in
Lemma~\ref{lem:affine-injectivity} is rational.
\end{proof}

\section{The exact mixed-adjoint input and forced scale}
For every integer $m\ge0$ and parameter $\gamma$, define
\[
 T_{m,\gamma}(t)=\sum_{k=0}^m\binom mk(\gamma)_{3k}t^k,\qquad
 B_{0,\gamma}=1,\quad
 B_{j+1,\gamma}=(z+3j+\gamma+2)B_{j,\gamma}
                       -zB_{j,\gamma}' .
\]
Thus $P_n=T_{n,2n+1}$. Set
\[
 M=\gamma+3m,\quad A=(M+3)(M+4),\quad
 \kappa=\frac{3(m+1)}{A},
\]
\[
 B=B_{m,\gamma},\qquad C=B_{m+1,\gamma},\qquad
 D_*=B_{m+2,\gamma},\qquad H=D_*+\kappa z^3B.
\]
For $u\ne0$, write the cyclic column of a polynomial $J$ as
\[
 J(z)=J_0(z^3)+zJ_1(z^3)+z^2J_2(z^3),\qquad
 {\bf j}(u)=(J_0(u),J_1(u),J_2(u))^{\mathsf T}.
\]
\begin{lemma}[Exact imported mixed-adjoint equivalence]\label{lem:mixed-input}
For $m\ge1$, $\gamma>0$ and $u\ne0$, the two equalities
\[
 T_{m,\gamma}(u^{-1})=T_{m+1,\gamma+2}(u^{-1})=0
\]
are equivalent to $\operatorname{rank}[{\bf c}(u),{\bf h}(u)]<2$.
The equivalence is algebraic in characteristic zero, so it also holds
over $K$ for positive rational $\gamma$.
\end{lemma}
\begin{proof}[Precise source of the imported identity]
This is Lemma \texttt{lem:mixed-input} of the
\href{QUINTIC_THREE_ADIC_PARITY_REDUCTION.tex}{three-adic parity note}.
Its original source is Theorem
\texttt{thm:gap-three-}\allowbreak\texttt{offset-two-rank} of the
\href{lower_exponent_initial_and_boundary.tex}{lower-exponent note}.
Those proofs specify the full Ore and concomitant identities,
include multiple lower roots, and justify their algebraic transfer
to $K$. Here only the necessary rank direction is used.
For the ordinary specialization $m=n$, $\gamma=2n+1$,
\[
 M=5n+1,\qquad
 \frac{3(n+1)}{(M+3)(M+4)}
 =\frac3{5(5n+4)}=\kappa .
\]
All auxiliary columns have the same fixed lower shape $\gamma$;
the upper moment has shape $\gamma+2$.
No uniform rank assertion or simplicity assumption is imported.
\end{proof}


\begin{lemma}[Every common root is balanced]\label{lem:balance}
If $P_n$ and $P_{n+1}$ have a common root $t\in K$, then
\[
 v_3(t)=1-r,\qquad
 \overline{t/3^{\,1-r}}=-\bigl(5\overline{D/3^r}\bigr)^{-1}\in k^\times.
\]
Equivalently $u=t^{-1}$ has valuation $r-1\ge1$ and
$\eta=\kappa u\in\mathcal O^\times$.
In particular a common root is impossible when $r\le1$.
\end{lemma}
\begin{proof}
Put $\gamma=2n+1$, $u=t^{-1}$ and choose $z^3=u$.
The constant terms are one, so $t,z,u$ are nonzero.
Every $(\gamma)_{3k}$ contains at least $k$ integer factors divisible
by three. Thus every coefficient of
\[
 z^{3n}P_n(z^{-3})
\]
at $z^{3(n-k)}$ has valuation at least $k$.
If $v_3(z)<1/3$, its leading term is uniquely of smallest valuation,
which is impossible at a root. Consequently $v_3(u)\ge1$.

Use Lemma~\ref{lem:mixed-input} with $m=n$:
\[
 B=B_{n,\gamma},\quad C=B_{n+1,\gamma},\quad
 D_*=B_{n+2,\gamma},\quad
 H=D_*+\kappa z^3B,\qquad
 \kappa=\frac3{5(5n+4)}.
\]
A common root forces the two cyclic columns $C,H$, evaluated at $u$,
to have rank less than two. Their coefficient vectors, as well as
that of $B$, are integral because $\gamma$ is an integer and
$v_3(u)>0$.
On reduction modulo the maximal ideal, $u$ is zero and the auxiliary
operator is
\[
 \mathsf A=z-zD_z+\bar\gamma-1
 \quad\hbox{on }k[z]/(z^3),\qquad \mathsf A^3=\mathsf A.
\]
Indeed its triangular diagonal is
$\bar\gamma-1,\bar\gamma-2,\bar\gamma-3$, the three elements
of $\mathbb F_3$, so its characteristic polynomial is $X^3-X$.
The vectors $1,\mathsf A1,\mathsf A^21$ form a basis since they are
monic of degrees zero, one and two. Moreover $\bar\gamma\in\mathbb F_3$.
For every $n\ge1$, $\mathsf A^n1$ and $\mathsf A^{n+1}1$
are independent: $\mathsf A$ has the three distinct eigenvalues
$0,1,-1$, the cyclic vector has a nonzero component in every
eigenspace, and the two nonzero eigenspaces give a Vandermonde minor.
The reduced columns $B,C,D_*$ are respectively
$\mathsf A^n1,\mathsf A^{n+1}1,\mathsf A^{n+2}1$.

Write $\eta=\kappa u$. If $v_3(\eta)>0$, the pair $C,H$ reduces
to $\mathsf A^{n+1}1,\mathsf A^{n+2}1$, of rank two.
If $v_3(\eta)<0$, first multiply the second column by $\eta^{-1}$;
it then reduces to the independent pair
$\mathsf A^{n+1}1,\mathsf A^n1$.
Both contradict rank deficiency. Hence $v_3(\eta)=0$.

Since $n\ge1$, $\mathsf A^{n+2}1=\mathsf A^n1$.
The pair now reduces to
$\mathsf A^{n+1}1,(1+\bar\eta)\mathsf A^n1$.
It has rank less than two exactly when $\bar\eta=-1$.
As $v_3(\kappa)=1-r$, balance gives $v_3(u)=r-1$
and $v_3(t)=1-r$.
Writing $t=3^{1-r}w$, one has
$\eta=1/[5(D/3^r)w]$, proving the residue assertion.
The inequality $v_3(u)\ge1$ excludes $r\le1$.
\end{proof}

\section{A scalar minor}
\begin{lemma}[One fixed row pair has affine reduction]
\label{lem:fixed-row-minor}
Suppose $r\ge2$ and put
\[
 q=\kappa^{-1}=\frac{5(5n+4)}3\in3\mathbb Z .
\]
With cyclic rows zero and one, the finite integer polynomial
\[
 F_n(X)=
 \det[{\bf c}(qX),{\bf d}(qX)+X{\bf b}(qX)]_{\{0,1\}}
                                                        \tag{2}\label{eq:scalar-minor}
\]
satisfies
\[
 \bar F_n(X)=\epsilon_n(1+X),\qquad
 \epsilon_n=\begin{cases}1,&n\text{ odd},\\-1,&n\text{ even}.
 \end{cases}
\]
Every common root of $P_n,P_{n+1}$ gives a zero
$X=\eta=\kappa/t$ of $F_n$ in $\mathcal O^\times$.
\end{lemma}
\begin{proof}
Here $r\ge2$ implies $n\equiv1\pmod9$ and $\gamma\equiv0\pmod3$.
On the reduced cyclic algebra,
\[
 A=z-zD_z-1,\qquad A1=z-1,\qquad A^21=z^2+1.
\]
The positive powers alternate between the last two vectors.
For odd $n$,
\[
 \bar{\bf b}(0)=(-1,1,0)^{\mathsf T},\qquad
 \bar{\bf c}(0)=(1,0,1)^{\mathsf T};
\]
for even $n$ these columns are reversed.
Always $\bar{\bf d}(0)=\bar{\bf b}(0)$.
The determinant of rows zero and one of
$[\bar{\bf c}(0),\bar{\bf b}(0)]$ is $\epsilon_n$.

All three cyclic columns belong to $\mathbb Z[u]^3$.
Since $q$ is divisible by three, substituting $u=qX$ reduces each
whole column polynomial to its value at zero.
This proves the asserted affine reduction, including every
coefficient of degree at least two.
The imported rank condition makes its fixed minor vanish at
$X=\kappa u$. Lemma~\ref{lem:balance} puts this point in
$\mathcal O^\times$.
\end{proof}

\begin{theorem}[All ordinary quintic consecutive moment pairs]
\label{thm:ordinary-quintic-coprime}
For every integer $n\ge1$,
\[
 \gcd(P_n,P_{n+1})=1
\]
over $\mathbb Q$ and over $\mathbb C$.
\end{theorem}
\begin{proof}
If $r\le1$, Lemma~\ref{lem:balance} excludes common roots over $K$.
If $r\ge2$, every common root gives a zero of $F_n$ in the integral
disk. Lemmas \ref{lem:affine-injectivity} and
\ref{lem:fixed-row-minor} show that at most one such zero exists.
The map $t\mapsto\kappa/t$ is injective on nonzero roots, so
$P_n,P_{n+1}$ have at most one distinct common root in $K$.

They have no common real root. One index is even, and for even $j$
and real $t$ finite Gamma-moment expansion gives
\[
 P_j(t)=\frac1{\Gamma(2j+1)}
       \int_0^\infty x^{2j}(1+tx^3)^j e^{-x}\,dx>0.
\]
The integrand is nonnegative and positive on an interval.
Apply Lemma~\ref{lem:real-field-descent} with $E=\mathbb Q$.
This concerns all conjugates of a rational common factor;
no simplicity assumption on the input polynomials is needed.
\end{proof}


\begin{corollary}[The ordinary support $\{2,5\}$]
\label{cor:quintic-sfc}
Every nonzero complex phase $X^2(A+BX^3)$ satisfies Strong Factorial:
its ordinary factorial moments cannot both vanish at any consecutive
positive orders.
\end{corollary}
\begin{proof}
For $AB\ne0$,
\[
 \mathcal L_1((X^2(A+BX^3))^j)=A^j(2j)!P_j(B/A).
\]
Theorem~\ref{thm:ordinary-quintic-coprime} excludes two consecutive
zeros. If $A=0$ or $B=0$, every positive-order factorial moment
is nonzero because the phase is a nonzero monomial.
\end{proof}

\begin{corollary}[All ordinary binomials of degree at most five]
\label{cor:degree-five-sfc}
Every nonzero ordinary univariate binomial of degree at most five
satisfies Strong Factorial at every consecutive positive pair.
\end{corollary}
\begin{proof}
Corollary \texttt{cor:degree-five-ray} of the
\href{QUINTIC_THREE_ADIC_PARITY_REDUCTION.tex}{three-adic parity note}
proves equivalence with the support $\{2,5\}$ tests at terminal
$N\ge11$, $N\equiv11\pmod{27}$.
Those tests are included in
Theorem~\ref{thm:ordinary-quintic-coprime}.
The support reduction and this theorem supply the entire argument;
no index-prefix experiment is used.
\end{proof}

\begin{corollary}[The exact ordinary degree-six frontier]
\label{cor:degree-six-frontier}
For
\[
 P^{a,d}_j(t)=\sum_{k=0}^j\binom jk(aj+1)_{dk}t^k,
\]
Strong Factorial for every ordinary univariate binomial of degree
at most six is equivalent to
\[
 \gcd_{\mathbb Q[t]}(P^{a,d}_{N-1},P^{a,d}_N)=1
\]
for every terminal $N\ge6$ and each of the three pairs
\[
 (a,d)=(2,4),\quad(3,3),\quad(4,2).
\]
These are the supports $\{2,6\}$, $\{3,6\}$ and $\{4,6\}$.
\end{corollary}
\begin{proof}
Corollary~\ref{cor:degree-five-sfc} settles all smaller degrees.
A nonmonomial binomial of degree six has the form
$X^a(A+BX^{6-a})$, where $0\le a\le5$ and $AB\ne0$.
For $a=0$, Theorem \texttt{thm:all} of the
\href{ALL_UNIVARIATE_TWO_TERM_SFC.tex}{constant-plus-power note}
applies. For $a=1$, gap five is covered by Corollary
\texttt{cor:prime-power-}\allowbreak\texttt{lower-one} of the
\href{ARITHMETIC_FRONTIER_TERMINAL_ORDERS.tex}{arithmetic note},
with prime five and ordinary shape one.
For $a=5$, Theorem \texttt{thm:main} of the
\href{UNIFORM_PADE_PARABOLA_SFC.tex}{adjacent-exponent note}
applies. This leaves exactly the three displayed pairs.

Theorem \texttt{thm:ordinary-fifth} of the
\href{lower_exponent_initial_and_boundary.tex}{lower-exponent note}
covers all ordinary binomials through terminal order five.
For the three remaining supports, the exact identity
\[
 \mathcal L_1((X^a(A+BX^d))^j)
       =A^j(aj)!P^{a,d}_j(B/A)
\]
identifies the remaining tests at every $N\ge6$ with the stated
rational gcds. Degenerate monomials have nonzero factorial moments.
Conversely, all the displayed tests are part of the degree-six
binomial assertion. This proves equivalence, without proving
any of the three unrestricted families.
\end{proof}

\paragraph{Replay and limits.}
The checker constructs eight full literal minors, verifies all
their coefficient reductions, and rejects a corrupted linear
coefficient. Complete reduced-operator cases verify the finite
linear algebra used above. The integrality and reduction proofs
supply the all-index conclusion.
This proves the stated ordinary quintic and degree-five binomial
claims; the degree-six corollary is a reduction to three families.
Unrestricted ordinary binomial SFC, general FC/SFC, and every real
Gamma shape are separate claims not proved here.
Internal replay does not establish formal verification, external
review or priority.
\end{document}
